\chapter{Laparoscopic inguinal hernia repair (TEP)}

\section*{Definition and Overview}
Laparoscopic Totally Extraperitoneal (TEP) inguinal hernia repair is a minimally invasive surgical procedure used to correct a groin hernia. The "Totally Extraperitoneal" approach means the surgeon accesses and repairs the hernia defect behind the abdominal wall muscles without entering the peritoneal (abdominal) cavity where the bowel and other organs reside. This technique uses small incisions, a laparoscope (a thin telescope with a light and camera), and a surgical mesh to reinforce the weakened groin area.

\section*{Epidemiology}
Inguinal hernia repair is one of the most common surgical procedures performed globally and in Australia. Inguinal hernias are far more common in males than females, with a lifetime risk of approximately 27\% for men and 3\% for women. TEP is particularly indicated for bilateral hernias, recurrent hernias (previously repaired via an open approach), or active patients seeking a quicker return to normal activities.

\section*{Aetiology and Pathophysiology}
Inguinal hernias occur when intra-abdominal pressure pushes tissue through a weakness in the inguinal canal. The TEP procedure addresses this mechanical defect through a physiological and anatomical approach:

\begin{itemize}
    \item \textbf{Anatomical weakness:} Natural defects in the fascia (transversalis fascia) or muscles of the groin canal create potential spaces for hernia development.
    \item \textbf{Raised intra-abdominal pressure:} Chronic coughing, heavy lifting, straining during bowel movements, or obesity can force abdominal contents through these weaknesses.
    \item \textbf{Extraperitoneal access:} By inflating the space between the rectus abdominis muscle and the peritoneum with carbon dioxide, the surgeon works in the preperitoneal space, avoiding contact with intra-abdominal organs.
    \item \textbf{Mesh reinforcement:} Placing a large synthetic mesh over the direct, indirect, and femoral hernia spaces (myopectineal orifice) redistributes abdominal pressure and prevents recurrence.
\end{itemize}

\section*{Clinical Features}
Patients undergoing or preparing for TEP inguinal hernia repair typically present with specific signs and symptoms:

\begin{itemize}
    \item \textbf{Groin lump:} A visible bulge in the groin area that may extend into the scrotum in men.
    \item \textbf{Discomfort or pain:} A dull ache, burning sensation, or dragging feeling in the groin, which often worsens with standing, lifting, or straining.
    \item \textbf{Reducibility:} The lump usually disappears when lying flat or can be gently pushed back.
    \item \textbf{Asymptomatic cases:} Some hernias are discovered incidentally during routine physical examinations.
\end{itemize}

\section*{Differential Diagnosis}
Several other conditions causing groin swelling or pain must be distinguished from an inguinal hernia:

\begin{itemize}
    \item \textbf{Femoral hernia:} A hernia occurring lower in the groin, through the femoral canal, more common in women.
    \item \textbf{Hydrocele or varicocele:} Fluid collection or enlarged veins within the scrotum.
    \item \textbf{Groin strain or sports hernia:} Musculoskeletal pain in the groin area without a true fascial defect.
    \item \textbf{Lymphadenopathy:} Swollen lymph nodes in the groin due to infection or malignancy.
    \item \textbf{Lipoma of the spermatic cord:} A benign fatty tumour mimicking a hernia bulge.
\end{itemize}

\section*{When to Seek Medical Care}
Patients with a known inguinal hernia should consult a surgeon to discuss repair options, particularly if the hernia is causing pain or growing. Immediate emergency evaluation is required if the hernia becomes irreducible, severely painful, or if systemic symptoms develop, as these indicate incarceration or strangulation.

\textbf{Call 000 (or local emergency services) for:}
\begin{itemize}
    \item A hernia bulge that cannot be pushed back (irreducible or incarcerated hernia)
    \item Severe, sudden, or worsening pain in the groin or abdomen
    \item Nausea, vomiting, or inability to pass wind or have a bowel movement
    \item Redness, warmth, or severe tenderness over the hernia bulge
    \item High fever, chills, or signs of severe systemic illness
\end{itemize}

\section*{Diagnosis and Investigations}
The diagnosis is primarily clinical, but investigations may be required to confirm or plan the surgery:

\begin{itemize}
    \item \textbf{Physical examination:} Standing and coughing during a clinical examination to make the hernia bulge prominent.
    \item \textbf{Ultrasound scan:} A high-resolution scan of the groin to confirm the hernia and distinguish it from other masses.
    \item \textbf{MRI or CT scan:} Used in complex cases or when the diagnosis remains unclear in patients with chronic groin pain.
    \item \textbf{Pre-operative assessment:} Routine blood tests, electrocardiogram (ECG), and general anaesthetic workup before surgery.
\end{itemize}

\section*{Management}
The management of inguinal hernias is primarily surgical, with TEP being a key minimally invasive technique:

\begin{itemize}
    \item \textbf{Surgical repair (TEP):} Accessing the preperitoneal space, reducing the hernia sac, and placing a mesh under general anaesthesia.
    \item \textbf{Post-operative pain management:} Standard paracetamol and anti-inflammatory medicines to control mild to moderate discomfort.
    \item \textbf{Early mobilisation:} Encouraging light walking immediately after surgery to reduce the risk of deep vein thrombosis and speed up recovery.
    \item \textbf{Activity modification:} Avoiding heavy lifting (greater than 5-10 kg) and strenuous exercise for 2 to 4 weeks.
    \item \textbf{Wound care:} Keeping the three small abdominal incision sites clean and dry, with waterproof dressings typically left on for several days.
\end{itemize}

\section*{Complications}
\begin{itemize}
    \item Chronic groin pain (post-herniorrhaphy neuralgia) due to nerve irritation or scarring
    \item Seroma or haematoma formation (fluid or blood collection in the groin space)
    \item Urinary retention, particularly common in older men post-operatively
    \item Hernia recurrence, though risk is low with mesh repair (typically less than 2-4\%)
    \item Infection of the surgical wounds or, rarely, the prosthetic mesh
    \item Injury to surrounding structures, such as the spermatic cord, testicle, bladder, or blood vessels
\end{itemize}

\section*{Prognosis and Outcomes}
The prognosis after TEP inguinal hernia repair is generally excellent. The minimally invasive nature of the procedure results in less post-operative pain, shorter hospital stays (often performed as a day-case surgery), and a faster return to normal activities compared to open repair. Most patients can return to light work and driving within a week, and full physical activities after 2 to 4 weeks. Recurrence rates are low when performed by an experienced laparoscopic surgeon.

\section*{Prevention and Self-Care}
\begin{itemize}
    \item Avoid heavy lifting where possible, and use proper lifting techniques when necessary
    \item Maintain a healthy weight to reduce pressure on the abdominal wall
    \item Prevent constipation by eating a high-fibre diet and drinking plenty of water
    \item Treat chronic coughs promptly with medical assistance
    \item Avoid smoking, which causes chronic coughing and weakens collagen synthesis in tissues
    \item Support the groin gently with a hand if coughing or sneezing during the early recovery period
\end{itemize}

\section*{Sources and Further Reading}
The following reputable organisations provide further information on laparoscopic inguinal hernia repair:

\begin{itemize}
    \item Royal Australasian College of Surgeons. \textit{Information on hernia types and surgical options.} https://www.surgeons.org/
    \item HerniaSurge. \textit{International guidelines for groin hernia management.} https://www.herniasurge.com/
    \item NHS. \textit{Patient guide to inguinal hernia repair, recovery, and risks.} https://www.nhs.uk/
    \item Mayo Clinic. \textit{Clinical overview of inguinal hernia diagnosis and surgical repair.} https://www.mayoclinic.org/
    \item Healthdirect Australia. \textit{Inguinal hernia information and advice for Australian patients.} https://www.healthdirect.gov.au/
\end{itemize}
